\documentclass[11pt]{amsart}
\usepackage{amsmath,amssymb,amsthm}
\usepackage[margin=1.1in]{geometry}
\usepackage{xcolor}
\newcommand{\R}{\mathbb{R}}
\newcommand{\fpl}[1]{(-\Delta_p)^{s}#1}
\newcommand{\al}{\alpha}
\title[Discussion note]{Discussion note: two-phase dead core for the
fractional $p$-Laplacian\\ \large status, one question about
\emph{Improved regularity for a nonlocal dead-core problem}, and a proposed argument}
\author{Hikmatullo Ismatov}
\date{\today}
\begin{document}
\maketitle

\section{Status in one paragraph}
For $\fpl{u}=u_+^{\gamma}-u_-^{\gamma}$ in $B_1$ ($0<\gamma<p-1$,
$p\ge2$), the target is growth of order $\al=\frac{sp}{p-1-\gamma}$ at
points with $u=|Du|=0$. The scaling exponent and the vanishing of the
right-hand side along blow-ups (via $(p-1)$-homogeneity, replacing
linearization) are done. Following Correa--dos Prazeres \S6, the
Liouville theorem is reformulated with a vanishing $1$-jet hypothesis,
avoiding affine subtraction; the jet-to-growth lemma transplants
verbatim. The interior $C^{1,\al_0}$ input is
Giovagnoli--Jesus--Silvestre (arXiv:2509.26565); their own remark
suggests the needed $f\in L^\infty$ extension is feasible for
$p<\frac{1}{1-s}$ (forcing $s>\frac12$, as in your paper). Two
caveats found: their constants blow up as $s\nearrow1$, so the local
two-phase theorem should be proved directly rather than by the
$s\nearrow1$ limit; and the tail bookkeeping below. Since we last spoke I
have checked the proposed (4.12) repair step by step (\S3, now sound in
outline) and intersected the constraint regions (\S4): the explicit
admissible region is nonempty, and the one genuinely binding constraint
turns out to be the size of the $C^{1,\al_0}$ exponent.

\section{The question about (4.12)}
In the proof of Theorem 4.1, the growth bound (4.11) for
$w_k=v_k(\cdot+h)-v_k(\cdot)$ rests on (4.8), which is justified for
$|x|\le\frac{1}{2r_k}$ (as (4.7) requires $r_kR<1/2$). Beyond that
radius, $u_k$ is only bounded, so $|w_k|\le 2(\theta_k(r_k)
r_k^{\frac{2s}{1-\gamma}})^{-1}$. Passing (4.12) to the limit via
[CS, ARMA 2011] requires $w_k\to w_0$ in $L^1(\R^n,\omega)$
\emph{globally} (their Remark~4.4 shows local convergence cannot
suffice), and the outer region contributes
$\sim\theta_k^{-1}r_k^{-\frac{2s\gamma}{1-\gamma}}$, which is bounded
only under a lower bound on $\theta_k$ that (4.5) does not seem to
provide. \textbf{Question: how is the far field handled in (4.12)?}
I may well be missing something standard.

\section{A proposed argument (checked step by step; sound in outline)}
Truncate: $\tilde w_k=w_k\chi_{B_{\rho_k}}$, $\rho_k=\frac{1}{4r_k}$.
Then $(-\Delta)^s\tilde w_k=f_k+H_k$ in $B_R$, where $H_k$ is the far
field as a datum. The point is that only the \emph{oscillation} of
$H_k$ matters:
\begin{enumerate}
\item $\|\nabla H_k\|_{L^\infty(B_R)}\lesssim
  C_h\rho_k^{\al-2s-2}+\theta_k^{-1}r_k^{2s+1-\al}\to0$, the second
  exponent being positive iff $\gamma<\frac{1}{1+2s}$ (the standing
  assumption). Hence $H_k=a_k+o(1)$, $a_k$ constant.
\item $a_k$ is bounded \emph{a posteriori}: $\tilde w_k$ has uniform
  local regularity and uniform $L^1_s$ norm, so
  $(-\Delta)^s\tilde w_k$ is uniformly bounded at a fixed point, while
  $f_k\to0$.
\item Stability now applies and gives
  $(-\Delta)^s(v_0(\cdot+h)-v_0)=a(h)$ with $a(h)$ finite.
\item The cocycle identity $w^{h_1+h_2}=w^{h_1}(\cdot+h_2)+w^{h_2}$
  forces $a(h)=b\cdot h$.
\item Second differences: $(-\Delta)^s\Delta_h^2v_0=0$ entire, with
  growth $(1+|x|)^{\al-2}$; Theorem~3.1 ($\nu=1$) makes $\Delta_h^2v_0$
  affine for every $h$, so $v_0$ is a polynomial of degree $\le2$ (by
  the growth $\al<1+2s$), and the vanishing $2$-jet at the branching
  point gives $v_0\equiv0$.
\end{enumerate}
No lower bound on $\theta_k$ is needed. The step I worried about most is
(2): it is \emph{not} circular --- the bound on $(-\Delta)^s\tilde w_k$ at
a fixed point splits into a local singular part (controlled by the uniform
interior $C^{2s+\gamma}$ estimate) and a tail (controlled using only the
growth bound $|w_k(y)|\le C_h|y|^{\al-1}$ on the truncated support, finite
iff $\gamma<\frac{1}{1+2s}$), neither of which uses $a_k$ or any bound on
$\theta_k$; $a_k$ is then solved for. The truncation is precisely what
manufactures the global $L^1_s$ convergence that (4.12) needs.

The same far-field-as-datum mechanism carries to $p\neq2$ for steps
(1)--(3): the increments solve the \emph{linear} equation
$L_{K_{v_k}}w=g$ with the kernel
$K_v\simeq|v(x)-v(y)|^{p-2}|x-y|^{-n-sp}$ of GJS \S7, and the oscillation
computation produces the additional condition $\gamma<\frac{p-1}{1+sp}$
(which coincides with $\frac{1}{1+2s}$ at $p=2$, stricter for $p>2$).
Steps (4)--(5), however, use translation invariance of $(-\Delta)^s$,
which the $v$-dependent kernel $K_v$ does \emph{not} have; so for $p\neq2$
I would drop the constant-$a(h)$ route and apply the vanishing-jet
Liouville to $v_0$ directly (where $a(h)$ never appears) --- this is one
of the points I would like to confirm with you.

\section{What I would like to discuss}
\begin{enumerate}
\item Whether the (4.12) argument of \S2--\S3 is right, or whether the
  far field has a standard justification I missed.
\item Whether, for $p\neq2$, the vanishing-jet Liouville on $v_0$ indeed
  removes the need for steps (4)--(5) (which rely on the translation
  invariance lost when $p\neq2$).
\item The one genuinely binding constraint: after intersecting the tail,
  jet-order and datum conditions (all explicit and jointly nonempty;
  $\nu=1$ alone suffices for $p>2$, the window staying open up to $s=1$),
  the binding requirement is $\al<1+\al_0$ with $\al_0$ the GJS exponent.
  Do we state the theorem with $\al_0$ as an abstract constant, or try to
  extract an explicit lower bound $\al_0(n,s,p)$ from their proof? Also:
  the kernel $K_v$ degenerates where $v$ is flat (worst at the branching
  point) --- is there a clean way to see $v_0$ is nondegenerate on an open
  set?
\item Proving the local two-phase $p$-dead-core theorem directly (the
  $s\nearrow1$ route is blocked by the non-uniformity of the GJS
  constants); it does not appear in the literature and may stand alone.
\item Whether to include the $f\in L^\infty$ extension of GJS for
  $p<\frac{1}{1-s}$ as a section of our paper.
\end{enumerate}

\medskip
\noindent\emph{All details, computations and sources are in the project
repository:} \texttt{github.com/hikmatulloismatov1999/fractional-p-dead-core}
\end{document}
